\section{서론}
\label{sec:introduction}

% P1: 맥락과 관측 수단
사람, 차량, 동물처럼 움직이는 다수의 에이전트는 대규모 시공간 궤적 데이터를 남긴다~\cite{Zheng15,Atluri18}.
이러한 궤적을 사건 및 환경 기록과 결합하면 시간에 따른 에이전트의 위치뿐 아니라 에이전트 사이의 상호작용과 환경의 변화를 함께 분석할 수 있다.
특히 활동 영역이 시간에 따라 변하고 생존 경쟁이 발생하는 환경에서는 어떤 에이전트가 살아남고 어떤 에이전트가 탈락하는지를 이해하기 위해 영역의 변화와 경쟁을 함께 고려해야 한다.
배틀로얄 게임은 이러한 환경을 대규모로 관측할 수 있게 해 준다.
PUBG: BATTLEGROUNDS(이하 PUBG)에서는 최대 100명의 플레이어가 1인, 2인 또는 4인 팀으로 하나의 맵에 투입되어 마지막 한 팀이 남을 때까지 경쟁한다.
경기가 진행되는 동안 안전구역은 정해진 단계에 따라 줄어들고, 그 밖에 있는 플레이어는 지속적인 피해를 입는다.
따라서 안전구역은 구역 피해 없이 활동할 수 있는 영역이며, 그 축소는 팀의 이동과 교전 조건을 변화시킨다.
PUBG는 약 10초 간격의 플레이어 위치 기록과 교전, 아이템, 구역 이벤트를 포함하는 경기 텔레메트리를 공개한다~\cite{PUBGAPI}.
맵과 모드별로 정해진 규칙의 경기가 반복되므로, 이 데이터는 줄어드는 영역 아래의 생존을 반복적으로 관찰할 수 있는 관측 환경이 된다.

% P2: 추정 과제와 활용 맥락
이러한 경기에서 안전구역의 단계 전환은 팀의 중요한 의사결정 시점이 된다.
새 단계가 시작되어 다음 안전구역이 공개되면, 각 팀은 구역이 줄어들기 전에 언제 어떤 경로로 이동할지, 즉 로테이션을 정하고 현재 위치를 지킬지 교전할지를 판단해야 한다.
이러한 판단을 분석할 때 필요한 양은 각 단계 경계에서 생존한 팀이 이어지는 한 단계 안에 완전히 탈락할 조건부 확률, 즉 다음 단계 탈락 위험이다.
이 양은 e스포츠에서도 활용될 수 있다.
프로 팀의 코치와 분석가에게는 단계별 로테이션 당시 팀이 처했던 위험을 복기하는 보조 지표가 될 수 있다.
여러 팀이 넓은 맵에 흩어져 동시에 움직이고 교전하는 대회 경기에서는 중계 옵저버가 관찰할 팀의 우선순위를 정하는 데 활용할 수 있다.
또한 개발사가 단계별 대기 시간, 축소 시간, 구역 밖 피해를 조정하려면 각 단계에서 관측되는 탈락 압박을 파악할 필요가 있다~\cite{PUBGBlueZone2026}.
본 연구는 공개 경기 기록을 활용한 관전자 관점의 위험 분석을 다루며, 특정 팀이 당시 알고 있던 정보만을 사용하는 의사결정 지원과는 구분된다.
반면 경기 단위의 우승확률은 경기 전체의 최종 결과를 요약하지만, 현재 생존한 팀이 이어지는 한 단계 안에 탈락할 확률을 직접 나타내지는 않는다.

% P3: 난점과 기존 설계의 한계
그러나 다음 단계 탈락 위험을 예측하기는 어렵다.
위험은 줄어드는 영역의 기하와 주변 팀과의 상호작용에 함께 달려 있고, 그 크기는 단계마다 달라진다.
또한 위치와 사건은 서로 다른 주기로 기록되므로 시간 정보를 정합적으로 맞춰야 한다.
기절한 플레이어와 탈락한 플레이어를 구분해야 하며, 단계 경계 이후의 정보가 입력에 포함되면 결과가 누설된다.
기존 PUBG 예측 연구는 생존과 경기 결과의 예측을 서로 다른 방식으로 다룬다.
Sagong 등~\cite{Sagong25a,Sagong25b}은 단계별 활동 점수를 이용하여 개별 플레이어의 이후 생존 여부를 예측한다.
KRAFTON의 예측기~\cite{KRAFTON}는 팀의 행동, 상태, 공간 정보를 팀 간 Transformer에 입력하고, Cox 계열의 학습 목적을 이용하여 미래 탈락 순서에 관한 상대위험 점수를 학습한 뒤 이를 우승확률로 변환한다.
그러나 위험집합 내 상대위험 점수와 정해진 다음 단계 안의 탈락확률은 서로 다른 추정 대상이다.
이러한 상대위험 학습 목적만으로는 팀 전체에 공통적인 단계별 위험 수준이 직접 식별되지 않으므로, 예측 구간과 단계별 기저위험을 명시적으로 다루는 설계와의 비교가 필요하다.
더욱이 두 기존 설계는 분석 단위, 데이터, 평가 지표가 모두 달라서, 보고된 성능을 직접 비교하는 것만으로는 어떤 설계 선택이 예측 성능을 좌우하는지 알기 어렵다.

% P4: 생태학적 관점과 적용
줄어드는 영역 안에서의 생존은 서식지 변화 아래에서 집단의 존속을 다루는 생태학적 문제와 맞닿아 있으며, 서식지 손실은 생물다양성에 크고 일관된 부정적 영향을 미친다~\cite{Fahrig03}.
본 연구는 생태학의 모형을 그대로 옮기는 대신, 환경 조건에 따른 집단의 존속과 기간별 소멸확률이라는 관점을 PUBG에 적용한다(\ref{sec:related:survival}절).
안전구역의 축소를 구역 피해 없이 활동할 수 있는 영역의 단계적 감소로, 팀의 탈락을 해당 집단의 존속 종료로 대응시킨다.
이를 바탕으로 다음 단계 탈락 위험을 단계를 시간 단위로 하는 조건부 이산시간 hazard로 정식화한다.
이동을 환경과 연결하는 관점~\cite{Nathan08}도 함께 참고하여, 각 팀의 상태를 구역 경계에 대한 노출, 다음 구역에 도달하기 위한 이동 부담, 주변의 경쟁 밀도, 남은 팀원 수로 구성한 서식지 맥락으로 표현한다.
이 설계는 단계별 탈락이라는 추정 대상에 맞추어 위험의 시간 단위와 환경적 맥락을 함께 구성한다.

% P5: 연구 질문과 접근
이 관점은 세 가지 질문을 제기한다.
\begin{itemize}
    \item \textbf{RQ1:}
    서식지 축소 관점에서 구성한 구역 경계에 대한 노출, 이동 부담, 주변 경쟁 밀도, 남은 팀원 수는 다음 단계 탈락 위험과 어떤 방향과 크기의 관련을 보이며, 그 관계는 단계와 코호트에 걸쳐 얼마나 일관되는가?

    \item \textbf{RQ2:}
    단계 조건부 hazard 목적과 서식지 맥락 표현을 결합한 설계는 동일한 데이터와 평가 조건에서 기존 설계의 적용판과 어떤 예측 성능 차이를 보이는가?

    \item \textbf{RQ3:}
    입력 표현, 학습 목적, 단계별 기저위험, 모델 구조를 각각 변경할 때 예측 성능은 얼마나 달라지며, 그 차이는 코호트에 따라 어떻게 달라지는가?
\end{itemize}

이 질문에 답하기 위해 맵, 시점, 팀 규모로 나눈 12개 코호트의 공개 경기 13,948개를 사용한다.
이를 위해 약 10초 간격의 위치 기록과 비주기적 사건 로그를 결합하여, 예측 시점 이후의 정보가 입력되지 않는 팀--단계 학습 표본을 구성하는 파이프라인을 구축한다.
이 파이프라인은 위치와 사건 기록을 공통 시간축에 정렬하고, 단계 경계 이후의 정보를 입력에서 배제하며, 탈락 라벨을 위치 스냅샷이 아닌 사건 기록으로 구성한다.
이어서 두 기존 설계를 동일한 팀--단계 표본에서 비교할 수 있도록 적용한다.
다른 조건을 고정한 상태에서 입력 표현, 학습 목적, 단계별 기저위험, 모델 구조를 한 번에 하나씩 바꾸어 각 요소의 조건부 기여를 평가한다.

% 연구 질문에 대응하는 기여
본 논문의 기여는 세 질문에 대한 답으로 정리된다.
\begin{itemize}
    \item \textbf{RQ1에 대한 답.}
    다음 단계 탈락 위험을 단계적 서식지 축소 관점에서 정식화하고, 서식지 맥락과 단계 조건부 이산시간 hazard를 결합한 해석 가능한 모델을 제시한다.
    이 모델에서 이동 부담과 주변 경쟁 밀도는 탈락 위험을 높이고 남은 팀원 수는 위험을 낮추며, 세 공변량의 추정된 효과 방향은 평가 가능한 모든 코호트에서 일관된다.
    다만 효과의 크기는 코호트마다 다르다.
    구역 경계 안쪽으로의 깊이가 갖는 보호 효과는 안전구역이 줄어들수록 커져 중반 단계 이후 모든 코호트에서 관측되는 반면, 주변 경쟁 밀도의 효과는 단계가 진행될수록 작아진다.

    \item \textbf{RQ2에 대한 답.}
    이 정식화에 맞춘 설계, 즉 단계 조건부 hazard 목적과 서식지 맥락 표현을 두 기존 설계의 적용판과 대비하여 제시한다.
    동일한 표본과 평가 조건에서 예측 성능을 비교한 결과, 본 연구의 설계는 평가한 모든 코호트에서 두 적용판보다 우수한 점추정치를 보인다.

    \item \textbf{RQ3에 대한 답.}
    통제된 비교를 통해 성능 차이에 대한 각 설계 요소의 조건부 기여를 평가한다.
    상대위험 기반 설계와의 비교에서는 학습 목적의 변경이, 활동 점수 설계와의 비교에서는 입력 표현의 변경이 주요 성능 차이를 설명하며, 후자의 개선은 주로 서식지 맥락을 포함할 때 관측된다.
    그래프와 시계열 구조의 추가 이득은 작고, 표본이 많은 일부 코호트에서만 관측된다.
\end{itemize}
